% Source: https://artofproblemsolving.com/community/c3h1721378_magic_circle_triangle
% Stable identity: aops-topic-1721378-q01
Let the common side sum be $T$ and let $c$ be the sum of the three corner numbers. Adding the three side sums counts every number once and the corners twice:
\[
3T=(1+2+\dots+9)+c=45+c .
\]
The corners are three different digits, so $1+2+3=6\le c\le7+8+9=24$, and therefore $17\le T\le23$.

Both extremes occur:

\begin{asy}
usepackage("amsmath");
size(8cm);
defaultpen(fontsize(10pt));
void tri(pair o,int[] c,int[] s){
pair[] P={o,o+(6,0),o+(3,5.196)};
pair[] Q={P[0]+(P[1]-P[0])/3,P[0]+2*(P[1]-P[0])/3,P[1]+(P[2]-P[1])/3,P[1]+2*(P[2]-P[1])/3,P[2]+(P[0]-P[2])/3,P[2]+2*(P[0]-P[2])/3};
draw(P[0]--P[1]--P[2]--cycle,gray+linewidth(0.6));
for(int i=0;i<3;++i){filldraw(circle(P[i],0.55),white,linewidth(0.8));label("$"+string(c[i])+"$",P[i]);}
for(int i=0;i<6;++i){filldraw(circle(Q[i],0.55),white,linewidth(0.8));label("$"+string(s[i])+"$",Q[i]);}
}
tri((0,0),new int[]{1,2,3},new int[]{5,9,4,8,6,7});
tri((9,0),new int[]{7,8,9},new int[]{2,6,1,5,3,4});
label("sum $17$",(3,-1.3));label("sum $23$",(12,-1.3));
label("\mbox{}",(min(currentpicture,user=true)+max(currentpicture,user=true))/2,black); // keeps the text after the figure black when LaTeX typesets the labels inline
\end{asy}

with corners $1,2,3$ and sides $1+5+9+2=2+4+8+3=3+6+7+1=17$, and with corners $7,8,9$ and sides $7+2+6+8=8+1+5+9=9+3+4+7=23$. The difference is $23-17=\boxed{6}$.
